\documentclass[10pt,a4paper]{amsart}
\usepackage[margin=19mm]{geometry}
\usepackage{amsmath,amssymb,mathtools}
\usepackage[scr=rsfs,cal=euler]{mathalfa}
\usepackage{xfrac}
\usepackage[unicode,hidelinks,bookmarks=false]{hyperref}
\newcommand{\ie}{i.e.\ }
\newcommand{\cf}{cf.\ }
\newcommand{\resp}{respectively\ }
\newcommand{\Subsection}{\subsection}
\providecommand{\nobreakdash}{\nobreak\hbox{-}}
\setlength{\emergencystretch}{2em}
\multlinegap0pt
\usepackage{bm}
\usepackage{amssymb}
\usepackage{tikz}
\usepackage{enumerate}
\newtheorem{lemma}{Lemma}
\newcommand{\curl}{{\rm curl\,}}
\newcommand{\dive}{{\rm div\,}}
\newcommand{\pa}{\partial}
\title{\bf elastic scattering by locally rough interfaces}
\author{Chengyu Wu, Yushan Xue, Jiaqing Yang}
\date{Source: arXiv:2511.18799v2 (CC BY 4.0)}
\begin{document}
\maketitle

\noindent\emph{Source and licence.} \bf elastic scattering by locally rough interfaces. Version arXiv:2511.18799v2 (CC BY 4.0), distributed by arXiv under the Creative Commons Attribution 4.0 International licence, http://creativecommons.org/licenses/by/4.0/, as recorded in the arXiv licence metadata for this exact version and shown on the abstract page https://arxiv.org/abs/2511.18799v2. Full licence notice: \texttt{licenses/arxiv-2511.18799-CC-BY-4.0.txt}.\par\smallskip

\noindent\textit{Editorial scope.} Counted unit: the article's lemma lem2.1 (source0.tex lines 203--226). The article's complete original proof of it is delivered with it (source0.tex lines 227--293, one \texttt{proof} environment of 67 lines). No mathematics is written, repaired or re-derived: the statement and the proof are the article's own lines, in the article's own notation, and no retained line is substituted or re-flowed.  No further source-local context is needed: every cross-reference of every retained line resolves inside the two delivered spans.  Nothing was omitted from any retained span, so the package carries no omission record.  No retained line contains a citation, so the wrapper carries no bibliography and no bibliography entry.  No retained line contains a figure environment, an image-inclusion command or drawing code, so no image is delivered and the package declares no asset dependency.  Editorial formatting only, in addition: an AMS \texttt{amsart} preamble that restates the article's own declarations and macros used by the retained lines, namely the environments \texttt{lemma} on separate counters, together with the standard packages those lines need.  The retained environments print in this document as Lemma 1 in order of appearance, whereas the article prints the same items with the numbering of its own full document.  The complete native source of the article as distributed in the arXiv e-print for this exact version is bundled unchanged as \texttt{sources/189677/source0.tex}.
\begin{lemma}\label{lem2.1}
	The following identities hold 
	\begin{align*}
	\mathbf{P}_{\tilde\mu,\tilde\lambda}\mathbf{u}=
	\left\{
	\begin{array}{ll}
		(\mu+\tilde{\mu})(\pa_{\boldsymbol{\tau}}\mathbf{u})^\perp+(2\mu+\lambda)\boldsymbol{\nu}\dive\mathbf{u}+\mu\boldsymbol{\tau}\dive^\perp\mathbf{u},~~~&{\rm if}~d=2,\\
		(\mu+\tilde{\mu})\mathcal{M}_{\boldsymbol{\nu}}(\mathbf{u})+(2\mu+\lambda)\boldsymbol{\nu}\dive\mathbf{u}-\mu\boldsymbol{\nu}\times\curl\mathbf{u},~~~&{\rm if}~d=3, 
	\end{array}
	\right.
	\end{align*}
	where $\mathcal{M}_{\boldsymbol{\nu}}$ is the differential operator defined by 
	\begin{align*}
	  \mathcal{M}_{\boldsymbol{\nu}}(\mathbf{u}):=
	  \left(
	  \begin{array}{l}
	  	\nu_2\pa_1u_2-\nu_1\pa_2u_2+\nu_3\pa_1u_3-\nu_1\pa_3u_3 \\
	  	\nu_1\pa_2u_1-\nu_2\pa_1u_1+\nu_3\pa_2u_3-\nu_2\pa_3u_3 \\
	  	\nu_1\pa_3u_1-\nu_3\pa_1u_1+\nu_2\pa_3u_2-\nu_3\pa_2u_2
	  \end{array}
	  \right)
	\end{align*}
	for $\mathbf{u}=(u_1,u_2,u_3)^\top$ and $\boldsymbol{\nu}=(\nu_1,\nu_2,\nu_3)^\top$. 
\end{lemma}

\begin{proof}
	From the definition of the generalized stress vector $	\mathbf{P}_{\tilde\mu,\tilde\lambda}\mathbf{u}$, it suffices to show that 
	\begin{align*}
	  \pa_{\boldsymbol{\nu}}\mathbf{u}=
	  \left\{
	  \begin{array}{ll}
	  	(\pa_{\boldsymbol{\tau}}\mathbf{u})^\perp+\boldsymbol{\nu}\dive\mathbf{u}+\boldsymbol{\tau}\dive^\perp\mathbf{u},~~~&{\rm if}~d=2,\\
	  	\mathcal{M}_{\boldsymbol{\nu}}(\mathbf{u})+\boldsymbol{\nu}\dive\mathbf{u}-\boldsymbol{\nu}\times\curl\mathbf{u},~~~&{\rm if}~d=3. 
	  \end{array}
	  \right.
	\end{align*}
	Direct calculation yields that 
	\begin{align*}
	  \begin{aligned}
	  &\;(\pa_{\boldsymbol{\tau}}\mathbf{u})^\perp+\boldsymbol{\nu}\dive\mathbf{u}+\boldsymbol{\tau}\dive^\perp\mathbf{u} \\
	   =&\left(
	  \begin{array}{l}
	  	\nu_2\pa_1u_2-\nu_1\pa_2u_2+\nu_1(\pa_1u_1+\pa_2u_2)+\nu_2(\pa_2u_1-\pa_1u_2) \\
	  	-\nu_2\pa_1u_1+\nu_1\pa_2u_1+\nu_2(\pa_1u_1+\pa_2u_2)-\nu_1(\pa_2u_1-\pa_1u_2)
	  \end{array}
	  \right) \\
	  =&\left(
	  \begin{array}{l}
	  	\nu_1\pa_1u_1+\nu_2\pa_2u_1 \\
	  	\nu_1\pa_1u_2+\nu_2\pa_2u_2
	  \end{array}
	  \right)
	  =\pa_{\boldsymbol{\nu}}\mathbf{u}, 
	  \end{aligned}
	\end{align*}
	and 
	\begin{align*}
	\begin{aligned}
	  &\;\mathcal{M}_{\boldsymbol{\nu}}(\mathbf{u})+\boldsymbol{\nu}\dive\mathbf{u}-\boldsymbol{\nu}\times\curl\mathbf{u} \\
	  =&\left(
	  \begin{array}{l}
	  	\nu_2\pa_1u_2-\nu_1\pa_2u_2+\nu_3\pa_1u_3-\nu_1\pa_3u_3 \\
	  	\nu_1\pa_2u_1-\nu_2\pa_1u_1+\nu_3\pa_2u_3-\nu_2\pa_3u_3 \\
	  	\nu_1\pa_3u_1-\nu_3\pa_1u_1+\nu_2\pa_3u_2-\nu_3\pa_2u_2
	  \end{array}
	  \right)+
	  \left(
	  \begin{array}{l}
	  	\nu_1(\pa_1u_1+\pa_2u_2+\pa_3u_3) \\
	  	\nu_2(\pa_1u_1+\pa_2u_2+\pa_3u_3) \\ 
	  	\nu_3(\pa_1u_1+\pa_2u_2+\pa_3u_3) 
	  \end{array}
	  \right) \\
	   &\;-\left(
	  \begin{array}{l}
	  	\nu_2(\pa_1u_2-\pa_2u_1)-\nu_3(\pa_3u_1-\pa_1u_3) \\
	  	\nu_3(\pa_2u_3-\pa_3u_2)-\nu_1(\pa_1u_2-\pa_2u_1) \\
	  	\nu_1(\pa_3u_1-\pa_1u_3)-\nu_2(\pa_2u_3-\pa_3u_2) 
	  \end{array}
	  \right) \\
	  =&\left(
	  \begin{array}{l}
	  	\nu_1\pa_1u_1+\nu_2\pa_2u_1+\nu_3\pa_3u_1 \\
	  	\nu_1\pa_1u_2+\nu_2\pa_2u_2+\nu_3\pa_3u_2  \\
	  	\nu_1\pa_1u_3+\nu_2\pa_2u_3+\nu_3\pa_3u_3 
	  \end{array}
	  \right) 
	  =\pa_{\boldsymbol{\nu}}\mathbf{u}, 
	  \end{aligned}
	\end{align*}
	which are the desired identities. 
\end{proof}

\end{document}
